\section{Results}
\label{sec:results}

We report what the telemetry supports, and we state where it does not support a claim. Every number below is a macro
generated from the stored outputs (Section~\ref{sec:methods}).

\subsection{Campaign and measurement validity}

All \NumSweeps{} sweeps have a complete telemetry row, and each row comes from the code path described in
Section~\ref{sec:methods}. Every sweep completed on its first attempt (\FirstPassFailures{} failures). In an earlier campaign with the same pass set, \texttt{opt} crashed on \ArchivedCrashes{} gsm kernels with \texttt{unify-loop-exits}. The pass
is replaced by \texttt{newgvn} in its R2 slot on every kernel, so no R2 sweep crashes. Table~\ref{tab:regimes} gives the sweeps
per regime.

\begin{table}[t]
\centering
\caption{Sweeps per regime.}
\label{tab:regimes}
\setlength{\tabcolsep}{3pt}
\begin{tabular}{|l|l|r|}
\hline
Regime & Purpose & Sweeps \\
\hline
R1 & Standard and size levels & 360 \\
\hline
R2 & Single passes & 5,040 \\
\hline
R3 & Ordered pairs & 13,500 \\
\hline
R4 & O3 prefixes & 3,900 \\
\hline
R5 & Pass interference & 2,400 \\
\hline
R6 & M\_pass sequences & 3,600 \\
\hline
Total &  & 28,800 \\
\hline
\end{tabular}
\end{table}


Benchmark outputs were compared with the canonical output of each kernel. Of the \OutputsChecked{} sweeps, \OutputsAgreeing{}
produced the canonical output. The \OutputsDiffering{} exceptions are sweeps of \texttt{mibench/\allowbreak{}susan-smoothing}, and each writes
an output image that differs from the canonical image in most of its bytes. Section~\ref{sec:threats} reports the investigation.
For each of the exceptions, the same optimized module lowered at \texttt{-O2} and interpreted with \texttt{lli} gives the canonical output, so the difference appears only in the \texttt{-O0} lowering. These comparisons are reproduced in the replication package. We have not identified the cause, and we report
these sweeps as exceptions rather than as matching outputs. For the bitcount kernel, the comparison was repeated under the
corrected masking rule, and \BitcountRecheckMatching{} of \BitcountRecheckChecked{} sweeps matched. Each kernel runs on one input,
so this agreement is evidence about those inputs only.

\subsection{Instruction-count interactions}
\label{sec:results-instcount}

The instruction count is exact and is measured before lowering, so these results do not depend on the timing noise or on the
lowering level. Figure~\ref{fig:instcount} shows the tensor of instruction-count interactions over the \NumThreePasses{} R3 passes,
averaged over the kernels. Of its \InstInteractionCells{} ordered pairs, \InstInteractionConsistentCells{} have a mean that
exceeds two standard errors across the kernels. The standard errors measure how much the kernels differ, because each
kernel's count is exact. Table~\ref{tab:instcount} lists the largest entries. The largest magnitude is
\InstInteractionMaxAbsPct{}\% of the canonical instruction count, and it belongs to a self-pair. In Figure~\ref{fig:instcount}, rows are the first pass and columns the second. Positive values mean that the combined change of a pair is larger than the sum of its single changes. The diagonal holds the self-interactions, which are left blank in the figure and reported in Table~\ref{tab:instcount}.

\PaperFigure{fig1_instcount_interaction_heatmap.pdf}{Instruction-count interaction of the \NumThreePasses{} R3 passes, averaged over kernels. \label{fig:instcount}}

\begin{table}[t]
\centering
\caption{Largest instruction-count interactions over kernels, as a percentage of the canonical instruction count, with standard errors across kernels. The standard errors measure how much the kernels differ; the counts are exact.}
\label{tab:instcount}
\setlength{\tabcolsep}{3pt}
\begin{tabular}{|l|l|r|r|l|}
\hline
First pass & Second pass & Mean (\%) & SE (\%) & Beyond 2 SE \\
\hline
loop-vectorize & loop-vectorize & -27.72 & 3.80 & yes \\
\hline
gvn & gvn & +16.57 & 0.82 & yes \\
\hline
early-cse & early-cse & +13.41 & 0.77 & yes \\
\hline
early-cse & gvn & +13.17 & 0.75 & yes \\
\hline
gvn & early-cse & +12.92 & 0.76 & yes \\
\hline
instcombine & instcombine & +9.11 & 0.80 & yes \\
\hline
simplifycfg & simplifycfg & +6.79 & 0.41 & yes \\
\hline
early-cse & instcombine & +5.96 & 0.45 & yes \\
\hline
gvn & instcombine & +5.51 & 0.47 & yes \\
\hline
instcombine & early-cse & +5.31 & 0.52 & yes \\
\hline
instcombine & gvn & +5.08 & 0.53 & yes \\
\hline
jump-threading & simplifycfg & +4.61 & 0.27 & yes \\
\hline
\end{tabular}
\end{table}


Several of the largest entries involve passes that remove redundant instructions. The two orders of \texttt{early-cse} and
\texttt{gvn} are both positive, which is what overlapping removals would produce: when the two passes remove the same redundant
instructions, the pair removes fewer than the sum of its single removals.

\textbf{Self-interactions.} The identity of Remark~\ref{rem:identity} holds wherever a second application leaves the bitcode
unchanged. Section~\ref{sec:results-canonical} counts \InstSelfIdentityCases{} such kernel-pass cases, and the identity holds
exactly in all of them, with \InstSelfIdentityMismatches{} mismatches. On those kernels, a positive self-interaction marks a
pass whose first application removes instructions, and a negative one marks a pass whose first application adds them. The sign
of a self-interaction is therefore not evidence of a second effect, and the self-pairs of \texttt{gvn}, \texttt{early-cse}, and
\texttt{loop-vectorize} in Table~\ref{tab:instcount} are read this way. Where the second application changes the IR again, the
identity does not hold, and the self-interaction mixes the second change with the first.

\subsection{Canonical form and repeated application}
\label{sec:results-canonical}

Table~\ref{tab:passchecks} reports the IR-level checks over \PassChecksKernels{} kernels. Each pass of the R3 set is applied to a
module, and then again to the result. The table compares two canonical forms as a controlled ablation of the canonicalization: the sweep canonical form, which applies \texttt{mem2reg}, \texttt{loop-simplify}, and \texttt{lcssa}, and the \texttt{mem2reg}-only form, which omits the \texttt{loop-simplify} and \texttt{lcssa} steps.

On the \texttt{mem2reg}-only form, every loop pass of the set changes every kernel, and the smallest count among the five loop
passes is \PassChecksLoopMinMemTwoReg{}. On the sweep canonical form, the smallest count is \PassChecksLoopMinCanonical{}, for
\texttt{loop-unroll}, which changes \PassChecksUnrollCanonicalChanged{} of \PassChecksKernels{} kernels there. So the canonical
form changes the effect of \texttt{loop-unroll} in particular. For every other pass, the two forms differ by at most
\PassChecksOtherMaxDiff{} kernels.

\begin{table}[t]
\centering
\caption{IR-level checks over 60 kernels. Changed counts the kernels where the pass changes the module, on the sweep canonical form and on the mem2reg-only form. Repeated counts the kernels where a second application changes it again, out of those where the first changed it.}
\label{tab:passchecks}
\setlength{\tabcolsep}{3pt}
\begin{tabular}{|l|r|r|r|r|}
\hline
Pass & Changed & Repeated & Mem2reg: changed & Mem2reg: repeated \\
\hline
instcombine & 60 & 30 & 60 & 30 \\
\hline
simplifycfg & 60 & 23 & 60 & 22 \\
\hline
early-cse & 60 & 53 & 59 & 46 \\
\hline
gvn & 60 & 27 & 60 & 20 \\
\hline
licm & 58 & 23 & 60 & 25 \\
\hline
loop-rotate & 60 & 32 & 60 & 31 \\
\hline
indvars & 59 & 42 & 60 & 42 \\
\hline
loop-unroll & 24 & 21 & 60 & 23 \\
\hline
loop-vectorize & 57 & 44 & 60 & 44 \\
\hline
slp-vectorizer & 31 & 19 & 32 & 20 \\
\hline
sccp & 60 & 26 & 60 & 26 \\
\hline
dce & 27 & 19 & 25 & 20 \\
\hline
dse & 24 & 19 & 22 & 20 \\
\hline
reassociate & 60 & 24 & 60 & 22 \\
\hline
jump-threading & 60 & 27 & 60 & 25 \\
\hline
\end{tabular}
\end{table}


Repeated application changes the IR again in \PassChecksRepeatMinPct{}\% to \PassChecksRepeatMaxPct{}\% of the kernels where the
first application changed it. So a pass that changes the module is not, in general, idempotent on these kernels. The identity of
Remark~\ref{rem:identity} is therefore used only where a second application is unchanged, and the \InstSelfIdentityCases{} such
cases are the ones on which it is tested.

\subsection{Timing interactions}
\label{sec:results-timing}

The timing measure is secondary, because its noise is large relative to the interactions. The median absolute paired delta of
the identical-IR sweeps is \NoiseMedianAbsPct{}\% of the canonical cost, and the ninety-fifth percentile is
\NoiseNinetyFifthAbsPct{}\%. Of the \InteractionCells{} entries of the normalized timing interaction (Equation~\ref{eq:timing}),
\InteractionSignificantCells{} exceed two standard errors (Table~\ref{tab:interaction}). Under a normal approximation, two
standard errors flag about five percent of entries by chance, so this count is larger than chance alone would give. The largest
magnitude is \InteractionMaxAbsPct{}\% of the canonical cost, which lies below the ninety-fifth percentile of the noise. The
entries that pass the test are therefore small, and we do not treat them as established interactions.

\begin{table}[t]
\centering
\caption{Largest pairwise interaction entries over kernels, as a percentage of the canonical cost, with standard errors. Each entry is the mean of $(\Delta_{ij}-\Delta_i-\Delta_j)/J_0$ over kernels.}
\label{tab:interaction}
\setlength{\tabcolsep}{3pt}
\begin{tabular}{|l|l|r|r|l|}
\hline
First pass & Second pass & Mean (\%) & SE (\%) & Beyond 2 SE \\
\hline
instcombine & jump-threading & +3.396 & 1.215 & yes \\
\hline
instcombine & instcombine & +3.274 & 1.385 & yes \\
\hline
gvn & instcombine & +3.061 & 1.362 & yes \\
\hline
jump-threading & loop-rotate & -2.962 & 0.743 & yes \\
\hline
simplifycfg & loop-rotate & -2.926 & 0.809 & yes \\
\hline
simplifycfg & instcombine & +2.901 & 1.270 & yes \\
\hline
loop-rotate & instcombine & +2.837 & 1.467 & no \\
\hline
loop-rotate & jump-threading & -2.789 & 0.876 & yes \\
\hline
jump-threading & instcombine & +2.679 & 1.276 & yes \\
\hline
gvn & loop-rotate & -2.625 & 0.745 & yes \\
\hline
instcombine & loop-rotate & +2.610 & 1.289 & yes \\
\hline
loop-rotate & simplifycfg & -2.591 & 0.815 & yes \\
\hline
\end{tabular}
\end{table}


The single-effect ratio baseline (Equation~\ref{eq:ratio}) puts most of its largest values on self-pairs, with standard
errors as large as their means, and none of its five largest entries passes the two-standard-error test (Table~\ref{tab:mpass}).
Table~\ref{tab:comparison} compares it with the canonical-cost normalized estimator on the same pairs.

\begin{table}[t]
\centering
\caption{Largest entries of the ratio estimator, $\mathbf{M}^{\mathrm{ratio}}_{\mathrm{pass}}$ (Equation~\ref{eq:ratio}), with standard errors. An entry is marked when its magnitude exceeds two standard errors.}
\label{tab:mpass}
\setlength{\tabcolsep}{3pt}
\begin{tabular}{|l|l|r|r|l|}
\hline
First pass & Second pass & Mean & SE & Beyond 2 SE \\
\hline
reassociate & reassociate & -3.7198 & 3.6405 & no \\
\hline
instcombine & loop-rotate & -2.7512 & 2.7181 & no \\
\hline
loop-vectorize & loop-vectorize & -2.2911 & 2.3976 & no \\
\hline
slp-vectorizer & slp-vectorizer & -1.2466 & 0.8554 & no \\
\hline
instcombine & dce & -0.9200 & 0.9568 & no \\
\hline
loop-rotate & gvn & -0.6911 & 0.2307 & yes \\
\hline
instcombine & early-cse & -0.6766 & 0.5147 & no \\
\hline
loop-rotate & jump-threading & -0.6730 & 0.1786 & yes \\
\hline
jump-threading & loop-rotate & -0.6389 & 0.1667 & yes \\
\hline
gvn & loop-rotate & -0.6141 & 0.2177 & yes \\
\hline
instcombine & dse & -0.5888 & 0.3660 & no \\
\hline
reassociate & licm & +0.5722 & 0.3749 & no \\
\hline
\end{tabular}
\end{table}


\begin{table}[t]
\centering
\caption{The same pass pairs under the ratio estimator (Equation~\ref{eq:ratio}) and under the timing interaction normalized by the canonical cost (Equation~\ref{eq:timing}). The ratio columns are dimensionless. The interaction columns are percentages of the canonical cost.}
\label{tab:comparison}
\setlength{\tabcolsep}{3pt}
\begin{tabular}{|l|l|r|r|r|r|}
\hline
First pass & Second pass & Ratio mean & Ratio SE & Mean (\%) & SE (\%) \\
\hline
dse & dse & +0.42 & 0.27 & +0.97 & 0.76 \\
\hline
licm & licm & +0.37 & 0.71 & +0.44 & 0.77 \\
\hline
early-cse & early-cse & +0.05 & 0.30 & +0.58 & 0.65 \\
\hline
gvn & gvn & +0.46 & 0.45 & -0.96 & 1.31 \\
\hline
sccp & sccp & -0.27 & 0.35 & -0.36 & 0.52 \\
\hline
jump-threading & loop-vectorize & -0.19 & 0.10 & -0.88 & 0.78 \\
\hline
\end{tabular}
\end{table}


The commutators are antisymmetric by construction (Table~\ref{tab:lie}). Of the \LieCells{} pairs, \LieSignificantCells{} differ
from zero by more than two standard errors, and the largest mean magnitude is \LieMaxAbs{} of the canonical cost, which lies
below the ninety-fifth percentile of the timing noise.

\begin{table}[t]
\centering
\caption{Largest mean Lie-bracket commutators over kernels, in units of the canonical cost, with standard errors.}
\label{tab:lie}
\setlength{\tabcolsep}{3pt}
\begin{tabular}{|l|l|r|r|l|}
\hline
First pass & Second pass & Mean & SE & Beyond 2 SE \\
\hline
loop-rotate & loop-vectorize & -0.0279 & 0.0131 & yes \\
\hline
loop-vectorize & loop-rotate & +0.0279 & 0.0131 & yes \\
\hline
early-cse & loop-rotate & +0.0153 & 0.0055 & yes \\
\hline
loop-rotate & early-cse & -0.0153 & 0.0055 & yes \\
\hline
licm & loop-rotate & -0.0147 & 0.0070 & yes \\
\hline
loop-rotate & licm & +0.0147 & 0.0070 & yes \\
\hline
simplifycfg & indvars & -0.0138 & 0.0092 & no \\
\hline
indvars & simplifycfg & +0.0138 & 0.0092 & no \\
\hline
gvn & licm & +0.0133 & 0.0061 & yes \\
\hline
licm & gvn & -0.0133 & 0.0061 & yes \\
\hline
instcombine & loop-vectorize & +0.0090 & 0.0070 & no \\
\hline
loop-vectorize & instcombine & -0.0090 & 0.0070 & no \\
\hline
\end{tabular}
\end{table}


\subsection{Cleanup sequences around loop-vectorize}

Table~\ref{tab:r5} gives the timing effect of each cleanup sequence on \texttt{loop-vectorize}. Of the \RFiveSequencesJudged{}
sequences, \RFiveSequencesSignificant{} have a mean effect beyond two standard errors. Single kernels do show effects larger than
their noise, and up to \RFiveKernelsBeyondNoiseMax{} of the \NumKernels{} kernels exceed two noise standard deviations for one
sequence. The identical-IR noise estimate covers binaries that are identical, so it does not cover the differences in code layout
and alignment that distinct binaries can show. We therefore read single-kernel effects of a few percent as unresolved.

\begin{table}[t]
\centering
\caption{Interference of cleanup sequences with \texttt{loop-vectorize}, as a percentage of the canonical cost. Mean, standard error, and median are over kernels. The last column counts kernels whose effect exceeds two noise standard deviations of that kernel.}
\label{tab:r5}
\setlength{\tabcolsep}{3pt}
\begin{tabular}{|l|l|r|r|r|r|}
\hline
Position & Sequence & Mean (\%) & SE (\%) & Median (\%) & Beyond noise \\
\hline
after & dce & +0.11 & 0.24 & -0.07 & 6 \\
\hline
after & instcombine & -2.75 & 1.13 & -0.98 & 23 \\
\hline
after & jump-threading & -0.35 & 0.98 & +0.23 & 28 \\
\hline
after & simplifycfg & -2.14 & 1.69 & -0.01 & 38 \\
\hline
before & instcombine & -1.79 & 0.65 & -0.77 & 24 \\
\hline
before & instcombine+jump-threading & -0.48 & 1.01 & +0.17 & 32 \\
\hline
before & jump-threading & -0.11 & 0.96 & -0.48 & 37 \\
\hline
before & jump-threading+simplifycfg & -1.92 & 1.56 & -0.21 & 32 \\
\hline
before & simplifycfg & -1.95 & 1.65 & -0.02 & 29 \\
\hline
before & simplifycfg+instcombine & -2.29 & 1.69 & -0.33 & 34 \\
\hline
before & simplifycfg+jump-threading & -2.38 & 1.63 & -0.30 & 35 \\
\hline
\end{tabular}
\end{table}


\subsection{Redundant passes and compile-time waste}
\label{sec:results-rpb}

A pass execution is redundant when the IR it leaves unchanged. Across all sweeps the mean redundant-pass ratio is
\RpbOverallMean{}. Table~\ref{tab:rpb} shows it per regime. The standard levels have the highest ratio, \RpbOneMean{}, and the
\texttt{-O3} prefixes have \RpbFourMean{}. Figure~\ref{fig:entropy} shows the ratio and the opcode entropy along the prefixes.
Charging each redundant execution with its share of its sweep's pass time (Equation~\ref{eq:waste}) gives
\WasteTotalSeconds{} seconds in total across all sweeps. That share is an assumption, because the telemetry records the pass time
of each sweep and not of each execution.

\PaperFigure{fig3_shannon_entropy_rpb_trajectory.pdf}{Mean opcode entropy (top) and mean redundant-pass ratio (bottom) along the stages of the \texttt{-O3} pipeline. \label{fig:entropy}}

\begin{table}[t]
\centering
\caption{Redundant pass events and estimated compile-time waste per regime. The ratio is the mean over sweeps of the fraction of pass executions that -print-changed reports as unchanged. Waste is the sum over sweeps of compile time multiplied by that ratio.}
\label{tab:rpb}
\setlength{\tabcolsep}{3pt}
\begin{tabular}{|l|l|r|r|}
\hline
Regime & Purpose & Mean ratio & Waste (s) \\
\hline
R1 & Standard and size levels & 0.844 & 0.36 \\
\hline
R2 & Single passes & 0.622 & 1.70 \\
\hline
R3 & Ordered pairs & 0.629 & 6.41 \\
\hline
R4 & O3 prefixes & 0.745 & 2.40 \\
\hline
R5 & Pass interference & 0.568 & 1.03 \\
\hline
R6 & M\_pass sequences & 0.559 & 1.49 \\
\hline
\end{tabular}
\end{table}


\subsection{Pareto fronts and hypervolume}

Across the \NumKernels{} kernels, the three-objective fronts contain \ParetoPoints{} non-dominated points. Table~\ref{tab:pareto}
gives the hypervolume per suite at the reference point of the canonical build, in units of $10^{-6}$. The mean over all kernels is
\MeanHypervolumeMicro{} in those units. These values are small because the reference point is the canonical build, and the
objective gains over it are small. Figure~\ref{fig:hv} shows how the hypervolume changes with the reference point. The compile reference is fixed, so the hypervolume is a reference-dependent summary. The speed
objective uses the timing, and the compile objective uses the optimizer wall time, so both carry the limits of
Section~\ref{sec:threats}.

\begin{table}[t]
\centering
\caption{Exact three-objective hypervolume at the reference point (0.5, 0.5, 0.5), in units of $10^{-6}$, and Pareto-front size, averaged over the kernels of each suite.}
\label{tab:pareto}
\setlength{\tabcolsep}{3pt}
\begin{tabular}{|l|r|r|r|}
\hline
Suite & Kernels & Mean hypervolume ($\times10^{-6}$) & Mean front size \\
\hline
mibench & 20 & 215.41 & 24.4 \\
\hline
polybench & 30 & 18.42 & 32.5 \\
\hline
tsvc\_coremark & 10 & 16.63 & 32.7 \\
\hline
\end{tabular}
\end{table}


\PaperFigure{fig4_pareto_3d_hypervolume_surface.pdf}{Mean exact hypervolume of the per-kernel fronts as the speedup and size reference points are lowered. \label{fig:hv}}

Of the \ParetoPoints{} front members, \ParetoMemberUnsatPct{}\% are UNSAT, \ParetoMemberNotComparablePct{}\% are not comparable,
and \ParetoMemberSatPct{}\% are SAT. The study therefore does not establish that every Pareto schedule refines its canonical
build, and the paper does not claim it.

\subsection{Refinement outcomes}
\label{sec:results-refinement}

Table~\ref{tab:smt} gives the block-local verdict of each sweep per regime. Over all sweeps, \SmtUnsatPct{}\% are UNSAT
(\SmtUnsatCount{} sweeps), \SmtSatPct{}\% are SAT (\SmtSatCount{}), and \SmtNotComparablePct{}\% are not comparable
(\SmtNotComparableCount{}). No sweep is UNKNOWN or ERROR, and every sweep's optimized IR was reproduced by re-running \texttt{opt}
(\SmtIrNotReproducedCount{} failures).

\begin{table}[t]
\centering
\caption{Block-local refinement outcome per sweep, as a percentage of the sweeps in each regime. Failed counts tool errors and optimized IR that did not reproduce.}
\label{tab:smt}
\setlength{\tabcolsep}{3pt}
\begin{tabular}{|l|r|r|r|r|r|r|}
\hline
Regime & Sweeps & UNSAT & SAT & Unknown & Not comparable & Failed \\
\hline
R1 & 360 & 31.7\% & 0.0\% & 0.0\% & 68.3\% & 0.0\% \\
\hline
R2 & 5,040 & 91.9\% & 3.3\% & 0.0\% & 4.8\% & 0.0\% \\
\hline
R3 & 13,500 & 61.3\% & 15.4\% & 0.0\% & 23.3\% & 0.0\% \\
\hline
R4 & 3,900 & 34.1\% & 0.5\% & 0.0\% & 65.4\% & 0.0\% \\
\hline
R5 & 2,400 & 39.3\% & 9.2\% & 0.0\% & 51.5\% & 0.0\% \\
\hline
R6 & 3,600 & 27.9\% & 14.2\% & 0.0\% & 57.8\% & 0.0\% \\
\hline
\end{tabular}
\end{table}


The verdicts are limited by coverage. Of the \BlocksMatched{} matched block pairs, \BlocksEncodablePct{}\% lie inside the fragment
of Section~\ref{sec:unsat}, and the rest contain an instruction the fragment does not model. Every not-comparable sweep has matched
labels, but none of its matched blocks lies in the fragment. Within the encodable blocks, \NamesRefine{} names refine and
\NamesViolated{} are violated. The share of not-comparable sweeps is highest in R1 and R4 and lowest in R2, as
Table~\ref{tab:smt} shows. The analysis does not classify which instructions cause the exclusion in each case.

Every SAT witness falls into exactly one of three categories, and the categories cover all \WitnessCatTotal{} SAT witnesses. Category
A is a flag gain: the optimized instruction differs from its canonical instruction only in wrap or exactness flags, and it carries
at least one flag that the canonical one lacks. Category B is a flag loss only. Category C is every other witness, which keeps its
text or changes it in some way other than flags. There are \WitnessCatA{} witnesses in A, \WitnessCatB{} in B, and \WitnessCatC{} in C.
Most SAT answers (\WitnessFlagOnly{} of \WitnessTotal{}, \WitnessFlagOnlyPct{}\%) name an instruction whose only change is in its
flags. Of the \SatSweeps{} SAT sweeps, \WitnessGainSweeps{} (\WitnessGainSweepsPct{}\%) contain a witness in category A. Across all
\WitnessTotal{} SAT witnesses, the optimizer added \WitnessGainedFlags{} flags: \texttt{nuw} \GainedNuw{} times, \texttt{nsw}
\GainedNsw{}, \texttt{nneg} \GainedNneg{}, and \texttt{samesign} \GainedSamesign{}. No witness gained an \texttt{exact} or
\texttt{disjoint} flag. The other witnesses changed the instruction in other ways (\WitnessOtherChange{}, of which
\WitnessRenamedOnly{} differ only in operand names), kept its text so that an upstream value explains the difference
(\WitnessUnchangedText{}), or lost a flag (\WitnessFlagLoss{}). A witness that only loses a flag cannot be a poison counterexample,
so its answer also depends on an upstream value.

Replaying each sweep one stage at a time locates the first stage that carries each gained flag on the same name. The replay
reproduces the optimized IR exactly in \ReplaysExact{} of \Replays{} replayed sweeps. Of the \IntroLocated{} located flag
introductions, \texttt{instcombine} accounts for \IntroInstcombine{} (\IntroInstcombinePct{}\%), \texttt{sccp} for \IntroSccp{}
(\IntroSccpPct{}\%), and \texttt{indvars} for \IntroIndvars{} (\IntroIndvarsPct{}\%). The other \IntroOther{} come from five passes.
For \IntroUntracked{} introductions, no stage kept the name, so they are not located. These counts identify the passes that add the
flags. They do not show why each flag is justified.

One justification is inspected in full. Sweep \texttt{K37-R2-018} runs \texttt{indvars} alone on the \texttt{patricia} kernel. Its
SAT witness is \texttt{\%157} in block \texttt{156} of \texttt{pat\_\allowbreak{}insert}, where the canonical instruction
\texttt{add nsw i32 \%.\allowbreak{}3, 1} becomes \texttt{add nuw nsw i32 \%.\allowbreak{}3, 1}. The counterexample is \texttt{\%.\allowbreak{}3 = -1}, for which the
unsigned increment wraps. The value \texttt{\%.\allowbreak{}3} is a loop counter, however. Its phi takes the value \texttt{1} on entry and
\texttt{\%157} on the back edge, and \texttt{\%157} is the only increment, so the counter never takes a negative value. The
block-local checker treats \texttt{\%.\allowbreak{}3} as a free input and admits \texttt{-1}, which the loop cannot produce. The same fact also
justifies the \texttt{samesign} flag that \texttt{indvars} gives the loop-header comparison, which contains a phi and so lies
outside the checker's fragment. This case is a false alarm of the block-local check, not a miscompilation. We have not inspected
the other sweeps, and this case does not show that they behave the same way.

A coarser count compares the number of flag tokens in each optimized module with the canonical module. By that count,
\SatAddingFlagPct{}\% of the SAT sweeps (\SatAddingFlagCount{} sweeps) gain a flag. We do not use it. It rises whenever a flagged
instruction appears anywhere in the module, including code that no witness names, so it cannot separate inference from copying.

Of the \OutputsChecked{} sweeps, \OutputsAgreeing{} matched the canonical output on the one input their kernel runs. None of the
\OutputsDiffering{} exceptions is a SAT sweep: all of them are UNSAT on their encodable blocks, so the block-local verdict does not
account for their output difference. The UNSAT verdicts are proofs only for the blocks that were compared, and the block-local
check does not cover floating-point or memory operations outside its fragment, so the verdicts do not show that the rest of each
program is unchanged. Figure~\ref{fig:z3} shows the validation time per sweep. It includes the re-derivation of the optimized IR
with \texttt{opt}, and it grows with the number of names compared.

Every refinement query of the Pareto-front sweeps is also exported as an \texttt{SMT-LIB 2} certificate. The export covers
\CertSweeps{} sweeps and \CertQueries{} queries, of which \CertSolverUnsat{} are UNSAT and \CertSolverSat{} are SAT. The z3
command-line binary, run independently on the exported text, returned the same verdict for \CertCliAgreement{} queries. The number of
sweeps whose status disagrees with the telemetry is \CertStatusDisagreements{}. The certificates are in the repository, one
gzip-compressed file per sweep.

\PaperFigure{fig5_z3_solver_scaling_profile.pdf}{Validation time per sweep against the number of names compared; sweeps with no compared name are plotted at one. \label{fig:z3}}
